\documentclass[10pt,a4paper]{amsart}
\usepackage[margin=19mm]{geometry}
\usepackage{amsmath,amssymb,mathtools}
\usepackage[scr=rsfs,cal=euler]{mathalfa}
\usepackage{xfrac}
\usepackage[unicode,hidelinks,bookmarks=false]{hyperref}
\newcommand{\ie}{i.e.\ }
\newcommand{\cf}{cf.\ }
\newcommand{\resp}{respectively\ }
\newcommand{\Subsection}{\subsection}
\providecommand{\nobreakdash}{\nobreak\hbox{-}}
\setlength{\emergencystretch}{2em}
\multlinegap0pt
\usepackage{mathtools}
\usepackage{bbm}
\usepackage{xcolor}
\usepackage{csquotes}
\usepackage{amssymb}
\usepackage{bm}
\usepackage{algorithm}\usepackage{algpseudocode}
\usepackage{caption}
\newtheorem{theorem}{Theorem}
\title{Affine Invariant Langevin Dynamics for rare-event sampling}
\author{Deepyaman Chakraborty, Ruben Harris, Rupert Klein, Guillermo Olic\'on-M\'endez, Sebastian Reich, Claudia Schillings}
\date{Source: arXiv:2601.00107v1 (CC BY 4.0)}
\begin{document}
\maketitle

\noindent\emph{Source and licence.} Affine Invariant Langevin Dynamics for rare-event sampling. Version arXiv:2601.00107v1 (CC BY 4.0), distributed by arXiv under the Creative Commons Attribution 4.0 International licence, http://creativecommons.org/licenses/by/4.0/, as recorded in the arXiv licence metadata for this exact version and shown on the abstract page https://arxiv.org/abs/2601.00107v1. Full licence notice: \texttt{licenses/arxiv-2601.00107-CC-BY-4.0.txt}.\par\smallskip

\noindent\textit{Editorial scope.} Counted unit: the article's theorem thm:consistency (source0.tex lines 369--384). The article's complete original proof of it is delivered with it (source0.tex lines 386--406, one \texttt{proof} environment of 21 lines). No mathematics is written, repaired or re-derived: the statement and the proof are the article's own lines, in the article's own notation, and no retained line is substituted or re-flowed.  Retained uncounted as the source-local context the proof needs: lines 351--359.  Nothing was omitted from any retained span, so the package carries no omission record.  No retained line contains a citation, so the wrapper carries no bibliography and no bibliography entry.  No retained line contains a figure environment, an image-inclusion command or drawing code, so no image is delivered and the package declares no asset dependency.  Editorial formatting only, in addition: an AMS \texttt{amsart} preamble that restates the article's own declarations and macros used by the retained lines, namely the environments \texttt{theorem} on separate counters, together with the standard packages those lines need.  The retained environments print in this document as Theorem 1 in order of appearance, whereas the article prints the same items with the numbering of its own full document.  The complete native source of the article as distributed in the arXiv e-print for this exact version is bundled unchanged as \texttt{sources/191969/source0.tex}.
\begin{equation}
    \label{eq:smooth_G}
\tilde G_\delta(x) :=
\begin{cases}
  0, & \text{if } G(x) < 0,\\[0.3em]
  G(x)\dfrac{\psi_\delta(G(x))}{\psi_\delta(G(x))+\psi_\delta(\delta-G(x))}, & \text{if } G(x)\in[0,\delta],\\[0.7em]
  G(x), & \text{if } G(x) > \delta.
\end{cases}
\end{equation}

\begin{theorem}[Consistency of the smoothed posterior]
\label{thm:consistency}
Let $F=\{x\in\mathbb X : G(x)\le 0\}$ be the failure domain and assume that $P_f=\mu_0(F)>0$. Let $\tilde G_\delta$ be the smoothed approximation defined in
\eqref{eq:smooth_G} and, for each $R>0$, define the probability measures
\begin{equation}\label{eq:smoothedpost}
  \mu_{\delta,R}(dx)
  := \frac{1}{Z_{\delta,R}}
     \exp\!\left(-\frac{\tilde G_\delta(x)^2}{2R}\right)\rho_0(x)\,dx,
  \quad  
  Z_{\delta,R}
  := \int_{\mathbb X} \exp\!\left(-\frac{\tilde G_\delta(x)^2}{2R}\right)
      \rho_0(x)\,dx.
\end{equation}
Then, for every fixed $\delta>0$ we have $\mu_{\delta,R}\;\xrightarrow{\mathrm{TV}}\;\mu_F$ as $R\rightarrow 0 $,
where the distribution $\mu_F(dx)=\frac{\mathbf{1}_F(x)}{\mu_0(F)}\rho_0(x)dx$ is independent of~$\delta$. 
\end{theorem}

\begin{proof}

Consider $\delta>0$ fixed. First, we show that the normalisation constants $Z_{\delta,R}$ converge to $\mu_0(F)$. Indeed, recall from \eqref{eq:smooth_G} that $\tilde{G}_\delta(x)=0$ whenever $x\in F$, for all $\delta\geq 0$. By splitting the state space $\mathbb{X}=F \cup (\mathbb X \setminus F)$, it follows from \eqref{eq:smoothedpost} that
\[
    Z_{\delta,R} = \mu_0(F) + \int_{\mathbb X \setminus F} \exp\left( -\frac{\tilde{G}_\delta(x^2)}{2R} \right) \rho_0(x)dx.
\]
The integrand on the right hand side is bounded by $\rho_0(x)$, and $\exp(-\tilde{G}(x)^2/(2R))\rightarrow 0$ for each $x\in \mathbb{X}\setminus F$. Hence, by the dominated convergence theorem we obtain that
\begin{equation}
    \label{eq:limit_normalisation}
        \lim_{R\rightarrow 0} Z_{\delta,R}=\mu_0(F).
\end{equation}

Denote $f_{\delta,R}(x):=\frac{\exp\left(- \tilde{G}(x)^2/Z_{\delta,R}\right)}{2R}\equiv \frac{N_{\delta,R}(x)}{Z_{\delta,R}}$ the density of $\mu_{\delta,R}$ with respect to $\mu_0$. We show that $f_{\delta,R}\rightarrow f_F(x):= \frac{1}{\mu_0(F)}\mathbf{1}_F(x)$ in $L^1(\mu_0)$. Indeed, by splitting again the state space and recalling that $\tilde G(x)=0$ for $x\in F$, we get
\begin{align*}
    \int_{\mathbb X}\vert f_{\delta,R}-f_F \vert d\mu_0 & = \int_F \left\vert \frac{N_{\delta,R}(x)}{Z_{\delta,R}} - \frac{1}{\mu_0(F)}\right\vert d\mu_0(x) + \int_{\mathbb X \setminus F} \frac{N_{\delta, R}(x)}{Z_{\delta,R}}d\mu_0(x)= \\
    &= \frac{\left\vert Z_{\delta,R}-\mu_0(F) \right\vert}{Z_{\delta,R}}+ \frac{1}{Z_{\delta,R}}\int_{\mathbb{X}\setminus F}N_{\delta,R}(x)d\mu_0(x).
\end{align*}
As $R\rightarrow 0$, the first term on the right hand side vanishes as given in \eqref{eq:limit_normalisation}. For the second term, observe that $N_{\delta,R}\rightarrow \mathbf{1}_F$ pointwise and that $N_{\delta,R}\leq 1 \in L^1(\mu_0)$, so that the integral on the right hand side vanishes as $R\rightarrow 0$ due to the dominated convergence theorem. 

Since the densities $f_{\delta,R}$ converge to $f_F$ in $L^1(\mu_0)$, their corresponding measures $\mu_{\delta,R}$ converge to $\mu_F$ in total variation as they are all absolutely continuous with respect to $\mu_0$, and the result follows.
\end{proof}

\end{document}
